\documentclass[12pt]{article}
\usepackage[utf8]{inputenc}
\usepackage[T1]{fontenc}
\usepackage{hyperref}
\title{Ulwa Language Resources for K-12 Curriculum Development}
\author{Compiled by Claude Sonnet 5 (claude-sonnet-5) for Andr\'as Kornai}
\date{}
\begin{document}
\maketitle

\section{Introduction}
Ulwa (ISO 639-3 ulw; Glottocode ulwa1239) is the southern variety of Sumu, a Misumalpan language of Nicaragua, centred on Karawala. It is endangered; one source reports only about 18 percent of those under 20 speak it. Miskitu is the dominant regional language.

\section{Primary Educational Resources}
Glottolog lists a 1995 CIDCA ``small dictionary'' with words and drawings, which may be usable as a school-level vocabulary book, but no copy or URL was found in the results. No textbooks or school programmes were found.

\section{Linguistic and Reference Materials}
\begin{itemize}
\item Glottolog languoid page and bibliography (includes the CIDCA 1989 and 1995 dictionaries, Koontz-Garboden 2009, Conzemius 1929): \url{https://glottolog.org/resource/languoid/id/ulwa1239}
\item NSF grant BCS-9308115 (project toward a reference grammar, dictionary and texts): \url{https://grantome.com/grant/NSF/BCS-9308115}
\item NSF grant BCS-9615545: \url{https://grantome.com/grant/NSF/BCS-9615545}
\item Article ``Ulwa (southern Sumo) - a distinct language?'' in \emph{Wani}: \url{https://camjol.info/index.php/WANI/article/view/19361}
\item Wani article PDF: \url{https://www.camjol.info/index.php/WANI/article/download/20000/24464/73553}
\item Hale, paper on the Ulwa research project: \url{https://ling-phil.mit.edu/papers/hale/papers/hale001.pdf}
\item ``Ulwa, Southern Sumu: the beginnings of a language research project'' (Cambridge, \emph{Linguistic Fieldwork}): \url{https://www.cambridge.org/core/books/abs/linguistic-fieldwork/ulwa-southern-sumu-the-beginnings-of-a-language-research-project/C62013CB3C6ADEC4F4A49FD7E03C39EF}
\item Numerals page for Ulwa, Karawala: \url{https://lingweb.eva.mpg.de/channumerals/Ulwa.htm}
\end{itemize}

\section{Religious and Community Resources}
A Cultural Survival article concerning the region was found: \url{https://www.culturalsurvival.org/index%2Ephp/node/13064}. No Ulwa scripture was found.

\section{Digital and Audio Resources}
None found.

\section{Curriculum Development Context}
Documentation exists mainly as academic and 1980s--90s CIDCA materials. A native-language curriculum would depend on locating the CIDCA dictionaries and the NSF-project outputs, and on working with Karawala community members. Details beyond these sources were not verified.
\end{document}